\documentclass[12pt]{article}

% Idioma y codificación
\usepackage[spanish, es-tabla]{babel}       %es-tabla para que se titule "Tabla"
\usepackage[utf8]{inputenc}

% Márgenes
\usepackage[a4paper,top=3cm,bottom=2.5cm,left=3cm,right=3cm]{geometry}

% Comentarios de bloque
\usepackage{verbatim}

% Paquetes de links
\usepackage[hidelinks]{hyperref}    % Permite enlaces
\usepackage{url}                    % redirecciona a la web

% Más opciones para enumeraciones
\usepackage{enumitem}

% Personalizar la portada
\usepackage{titling}


% Paquetes de tablas
\usepackage{multirow}


%------------------------------------------------------------------------

%Paquetes de figuras
\usepackage{caption}
\usepackage{subcaption} % Figuras al lado de otras
\usepackage{float}      % Poner figuras en el sitio indicado H.


% Paquetes de imágenes
\usepackage{graphicx}       % Paquete para añadir imágenes
\usepackage{transparent}    % Para manejar la opacidad de las figuras

% Paquete para usar colores
\usepackage[dvipsnames]{xcolor}
\usepackage{pagecolor}      % Para cambiar el color de la página

% Habilita tamaños de fuente mayores
\usepackage{fix-cm}


%------------------------------------------------------------------------

% Paquetes de matemáticas
\usepackage{mathtools, amsfonts, amssymb, mathrsfs}
\usepackage[makeroom]{cancel}     % Simplificar tachando
\usepackage{polynom}    % Divisiones y Ruffini
\usepackage{units} % Para poner fracciones diagonales con \nicefrac

\usepackage{pgfplots}   %Representar funciones
\pgfplotsset{compat=1.18}  % Versión 1.18

\usepackage{tikz-cd}    % Para usar diagramas de composiciones
\usetikzlibrary{calc}   % Para usar cálculo de coordenadas en tikz

%Definición de teoremas, etc.
\usepackage{amsthm}
%\swapnumbers   % Intercambia la posición del texto y de la numeración

\theoremstyle{plain}

\makeatletter
\@ifclassloaded{article}{
  \newtheorem{teo}{Teorema}[section]
}{
  \newtheorem{teo}{Teorema}[chapter]  % Se resetea en cada chapter
}
\makeatother

\newtheorem{coro}{Corolario}[teo]           % Se resetea en cada teorema
\newtheorem{prop}[teo]{Proposición}         % Usa el mismo contador que teorema
\newtheorem{lema}[teo]{Lema}                % Usa el mismo contador que teorema

\theoremstyle{remark}
\newtheorem*{observacion}{Observación}

\theoremstyle{definition}

\makeatletter
\@ifclassloaded{article}{
  \newtheorem{definicion}{Definición} [section]     % Se resetea en cada chapter
}{
  \newtheorem{definicion}{Definición} [chapter]     % Se resetea en cada chapter
}
\makeatother

\newtheorem*{notacion}{Notación}
\newtheorem*{ejemplo}{Ejemplo}
\newtheorem*{ejercicio*}{Ejercicio}             % No numerado
\newtheorem{ejercicio}{Ejercicio} [section]     % Se resetea en cada section


% Modificar el formato de la numeración del teorema "ejercicio"
\renewcommand{\theejercicio}{%
  \ifnum\value{section}=0 % Si no se ha iniciado ninguna sección
    \arabic{ejercicio}% Solo mostrar el número de ejercicio
  \else
    \thesection.\arabic{ejercicio}% Mostrar número de sección y número de ejercicio
  \fi
}


% \renewcommand\qedsymbol{$\blacksquare$}         % Cambiar símbolo QED
%------------------------------------------------------------------------

% Paquetes para encabezados
\usepackage{fancyhdr}
\pagestyle{fancy}
\fancyhf{}

\newcommand{\helv}{ % Modificación tamaño de letra
\fontfamily{}\fontsize{12}{12}\selectfont}
\setlength{\headheight}{15pt} % Amplía el tamaño del índice


%\usepackage{lastpage}   % Referenciar última pag   \pageref{LastPage}
\fancyfoot[C]{\thepage}

%------------------------------------------------------------------------

% Conseguir que no ponga "Capítulo 1". Sino solo "1."
\makeatletter
\@ifclassloaded{book}{
  \renewcommand{\chaptermark}[1]{\markboth{\thechapter.\ #1}{}} % En el encabezado
    
  \renewcommand{\@makechapterhead}[1]{%
  \vspace*{50\p@}%
  {\parindent \z@ \raggedright \normalfont
    \ifnum \c@secnumdepth >\m@ne
      \huge\bfseries \thechapter.\hspace{1em}\ignorespaces
    \fi
    \interlinepenalty\@M
    \Huge \bfseries #1\par\nobreak
    \vskip 40\p@
  }}
}
\makeatother

%------------------------------------------------------------------------
% Paquetes de cógido
\usepackage{minted}
\renewcommand\listingscaption{Código fuente}

\usepackage{fancyvrb}
% Personaliza el tamaño de los números de línea
\renewcommand{\theFancyVerbLine}{\small\arabic{FancyVerbLine}}

% Estilo para C++
\newminted{cpp}{
    frame=lines,
    framesep=2mm,
    baselinestretch=1.2,
    linenos,
    escapeinside=||
}



\usepackage{listings} % Para incluir código desde un archivo

\renewcommand\lstlistingname{Código Fuente}
\renewcommand\lstlistlistingname{Índice de Códigos Fuente}

% Definir colores
\definecolor{vscodepurple}{rgb}{0.5,0,0.5}
\definecolor{vscodeblue}{rgb}{0,0,0.8}
\definecolor{vscodegreen}{rgb}{0,0.5,0}
\definecolor{vscodegray}{rgb}{0.5,0.5,0.5}
\definecolor{vscodebackground}{rgb}{0.97,0.97,0.97}
\definecolor{vscodelightgray}{rgb}{0.9,0.9,0.9}

% Configuración para el estilo de C similar a VSCode
\lstdefinestyle{vscode_C}{
  backgroundcolor=\color{vscodebackground},
  commentstyle=\color{vscodegreen},
  keywordstyle=\color{vscodeblue},
  numberstyle=\tiny\color{vscodegray},
  stringstyle=\color{vscodepurple},
  basicstyle=\scriptsize\ttfamily,
  breakatwhitespace=false,
  breaklines=true,
  captionpos=b,
  keepspaces=true,
  numbers=left,
  numbersep=5pt,
  showspaces=false,
  showstringspaces=false,
  showtabs=false,
  tabsize=2,
  frame=tb,
  framerule=0pt,
  aboveskip=10pt,
  belowskip=10pt,
  xleftmargin=10pt,
  xrightmargin=10pt,
  framexleftmargin=10pt,
  framexrightmargin=10pt,
  framesep=0pt,
  rulecolor=\color{vscodelightgray},
  backgroundcolor=\color{vscodebackground},
}

%------------------------------------------------------------------------

% Comandos definidos
\newcommand{\bb}[1]{\mathbb{#1}}
\newcommand{\cc}[1]{\mathcal{#1}}

% I prefer the slanted \leq
\let\oldleq\leq % save them in case they're every wanted
\let\oldgeq\geq
\renewcommand{\leq}{\leqslant}
\renewcommand{\geq}{\geqslant}

% Si y solo si
\newcommand{\sii}{\iff}

% Letras griegas
\newcommand{\eps}{\epsilon}
\newcommand{\veps}{\varepsilon}
\newcommand{\lm}{\lambda}

\newcommand{\ol}{\overline}
\newcommand{\ul}{\underline}
\newcommand{\wt}{\widetilde}
\newcommand{\wh}{\widehat}

\let\oldvec\vec
\renewcommand{\vec}{\overrightarrow}

% Derivadas parciales
\newcommand{\del}[2]{\frac{\partial #1}{\partial #2}}
\newcommand{\Del}[3]{\frac{\partial^{#1} #2}{\partial^{#1} #3}}
\newcommand{\deld}[2]{\dfrac{\partial #1}{\partial #2}}
\newcommand{\Deld}[3]{\dfrac{\partial^{#1} #2}{\partial^{#1} #3}}


\newcommand{\AstIg}{\stackrel{(\ast)}{=}}
\newcommand{\Hop}{\stackrel{L'H\hat{o}pital}{=}}

\newcommand{\red}[1]{{\color{red}#1}} % Para integrales, destacar los cambios.

% Método de integración
\newcommand{\MetInt}[2]{
    \left[\begin{array}{c}
        #1 \\ #2
    \end{array}\right]
}

% Declarar aplicaciones
% 1. Nombre aplicación
% 2. Dominio
% 3. Codominio
% 4. Variable
% 5. Imagen de la variable
\newcommand{\Func}[5]{
    \begin{equation*}
        \begin{array}{rrll}
            #1:& #2 & \longrightarrow & #3\\
               & #4 & \longmapsto & #5
        \end{array}
    \end{equation*}
}

%------------------------------------------------------------------------
\let\oldRe\Re % save them in case they're every wanted
\let\oldIm\Im
\renewcommand{\Re}{\operatorname{Re}} % redefine them
\renewcommand{\Im}{\operatorname{Im}}
\DeclareMathOperator{\Log}{Log}
\DeclareMathOperator{\Arg}{Arg}

\begin{document}

    % 1. Foto de fondo
    % 2. Título
    % 3. Encabezado Izquierdo
    % 4. Color de fondo
    % 5. Coord x del titulo
    % 6. Coord y del titulo
    % 7. Fecha

    
    % 1. Foto de fondo
% 2. Título
% 3. Encabezado Izquierdo
% 4. Color de fondo
% 5. Coord x del titulo
% 6. Coord y del titulo
% 7. Fecha

\newcommand{\portada}[7]{

    \portadaBase{#1}{#2}{#3}{#4}{#5}{#6}{#7}
    \portadaBook{#1}{#2}{#3}{#4}{#5}{#6}{#7}
}

\newcommand{\portadaExamen}[7]{

    \portadaBase{#1}{#2}{#3}{#4}{#5}{#6}{#7}
    \portadaArticle{#1}{#2}{#3}{#4}{#5}{#6}{#7}
}




\newcommand{\portadaBase}[7]{

    % Tiene la portada principal y la licencia Creative Commons
    
    % 1. Foto de fondo
    % 2. Título
    % 3. Encabezado Izquierdo
    % 4. Color de fondo
    % 5. Coord x del titulo
    % 6. Coord y del titulo
    % 7. Fecha
    
    
    \thispagestyle{empty}               % Sin encabezado ni pie de página
    \newgeometry{margin=0cm}        % Márgenes nulos para la primera página
    
    
    % Encabezado
    \fancyhead[L]{\helv #3}
    \fancyhead[R]{\helv \nouppercase{\leftmark}}
    
    
    \pagecolor{#4}        % Color de fondo para la portada
    
    \begin{figure}[p]
        \centering
        \transparent{0.3}           % Opacidad del 30% para la imagen
        
        \includegraphics[width=\paperwidth, keepaspectratio]{assets/#1}
    
        \begin{tikzpicture}[remember picture, overlay]
            \node[anchor=north west, text=white, opacity=1, font=\fontsize{60}{90}\selectfont\bfseries\sffamily, align=left] at (#5, #6) {#2};
            
            \node[anchor=south east, text=white, opacity=1, font=\fontsize{12}{18}\selectfont\sffamily, align=right] at (9.7, 3) {\textbf{\href{https://losdeldgiim.github.io/}{Los Del DGIIM}}};
            
            \node[anchor=south east, text=white, opacity=1, font=\fontsize{12}{15}\selectfont\sffamily, align=right] at (9.7, 1.8) {Doble Grado en Ingeniería Informática y Matemáticas\\Universidad de Granada};
        \end{tikzpicture}
    \end{figure}
    
    
    \restoregeometry        % Restaurar márgenes normales para las páginas subsiguientes
    \pagecolor{white}       % Restaurar el color de página
    
    
    \newpage
    \thispagestyle{empty}               % Sin encabezado ni pie de página
    \begin{tikzpicture}[remember picture, overlay]
        \node[anchor=south west, inner sep=3cm] at (current page.south west) {
            \begin{minipage}{0.5\paperwidth}
                \href{https://creativecommons.org/licenses/by-nc-nd/4.0/}{
                    \includegraphics[height=2cm]{assets/Licencia.png}
                }\vspace{1cm}\\
                Esta obra está bajo una
                \href{https://creativecommons.org/licenses/by-nc-nd/4.0/}{
                    Licencia Creative Commons Atribución-NoComercial-SinDerivadas 4.0 Internacional (CC BY-NC-ND 4.0).
                }\\
    
                Eres libre de compartir y redistribuir el contenido de esta obra en cualquier medio o formato, siempre y cuando des el crédito adecuado a los autores originales y no persigas fines comerciales. 
            \end{minipage}
        };
    \end{tikzpicture}
    
    
    
    % 1. Foto de fondo
    % 2. Título
    % 3. Encabezado Izquierdo
    % 4. Color de fondo
    % 5. Coord x del titulo
    % 6. Coord y del titulo
    % 7. Fecha


}


\newcommand{\portadaBook}[7]{

    % 1. Foto de fondo
    % 2. Título
    % 3. Encabezado Izquierdo
    % 4. Color de fondo
    % 5. Coord x del titulo
    % 6. Coord y del titulo
    % 7. Fecha

    % Personaliza el formato del título
    \pretitle{\begin{center}\bfseries\fontsize{42}{56}\selectfont}
    \posttitle{\par\end{center}\vspace{2em}}
    
    % Personaliza el formato del autor
    \preauthor{\begin{center}\Large}
    \postauthor{\par\end{center}\vfill}
    
    % Personaliza el formato de la fecha
    \predate{\begin{center}\huge}
    \postdate{\par\end{center}\vspace{2em}}
    
    \title{#2}
    \author{\href{https://losdeldgiim.github.io/}{Los Del DGIIM}}
    \date{Granada, #7}
    \maketitle
    
    \tableofcontents
}




\newcommand{\portadaArticle}[7]{

    % 1. Foto de fondo
    % 2. Título
    % 3. Encabezado Izquierdo
    % 4. Color de fondo
    % 5. Coord x del titulo
    % 6. Coord y del titulo
    % 7. Fecha

    % Personaliza el formato del título
    \pretitle{\begin{center}\bfseries\fontsize{42}{56}\selectfont}
    \posttitle{\par\end{center}\vspace{2em}}
    
    % Personaliza el formato del autor
    \preauthor{\begin{center}\Large}
    \postauthor{\par\end{center}\vspace{3em}}
    
    % Personaliza el formato de la fecha
    \predate{\begin{center}\huge}
    \postdate{\par\end{center}\vspace{5em}}
    
    \title{#2}
    \author{\href{https://losdeldgiim.github.io/}{Los Del DGIIM}}
    \date{Granada, #7}
    \thispagestyle{empty}               % Sin encabezado ni pie de página
    \maketitle
    \vfill
}
    \portadaExamen{ffccA4.jpg}{Variable Compleja I\\Examen IV}{Variable Compleja I. Examen IV}{MidnightBlue}{-9.5}{28}{2024-2025}{Arturo Olivares Martos}

    \begin{description}
        \item[Asignatura] Variable Compleja I.
        \item[Curso Académico] 2021-22.
        \item[Grado] Doble Grado en Ingeniería Informática y Matemáticas.
        \item[Grupo] Único.
        \item[Profesor] Javier Merí de la Maza.
        \item[Descripción] Prueba Intermedia.
        \item[Fecha] 4 de Mayo de 2022.
        \item[Duración] 120 minutos.
    
    \end{description}
    \newpage

    \begin{ejercicio}[4 puntos]
        Probar que la serie $\displaystyle \sum_{n=0}^{\infty} \frac{\cos(nz)}{3^n}$ converge absolutamente en todo punto del dominio $\Omega = \{ z \in \mathbb{C} : -\ln(3) < \Im z < \ln(3) \}$. Estudiar la convergencia uniforme en subconjuntos de $\Omega$. Probar que la función $g : \Omega \to \mathbb{C}$ dada por
        \[
            g(z) = \sum_{n=0}^{\infty} \frac{\cos(nz)}{3^n}
        \]
        es continua en $\Omega$ y calcular $\displaystyle \int_{C(0,1)} \dfrac{g(z)}{z} \, dz$.
    \end{ejercicio}

    \begin{ejercicio}[3 puntos]
        Estudiar la derivabilidad de las funciones $f, g : \mathbb{C} \to \mathbb{C}$ dadas por
        \[
            f(z) = \sen\left(\ol{z}\right)\qquad g(z) = z(z-1)f(z) \quad (z \in \mathbb{C}).
        \]
    \end{ejercicio}

    \begin{ejercicio}[3 puntos]
        Sean $a, b \in \mathbb{C}$ con $a \neq b$ y sea $R > 0$ de modo que $R > \max\{ |a|, |b| \}$. Probar que, si $f$ es una función entera, se tiene que:
        \[
            \int_{C(0,R)} \frac{f(z)}{(z-a)(z-b)} \, dz = 2\pi i \cdot \frac{f(b) - f(a)}{b-a}.
        \]
        Deducir que toda función entera y acotada es constante (Teorema de Liouville).
    \end{ejercicio}


    \newpage
    \setcounter{ejercicio}{0}

    \begin{ejercicio}[4 puntos]
        Probar que la serie $\displaystyle \sum_{n=0}^{\infty} \frac{\cos(nz)}{3^n}$ converge absolutamente en todo punto del dominio $\Omega = \{ z \in \mathbb{C} : -\ln(3) < \Im z < \ln(3) \}$. Estudiar la convergencia uniforme en subconjuntos de $\Omega$. Probar que la función $g : \Omega \to \mathbb{C}$ dada por
        \[
            g(z) = \sum_{n=0}^{\infty} \frac{\cos(nz)}{3^n}
        \]
        es continua en $\Omega$ y calcular $\displaystyle \int_{C(0,1)} \dfrac{g(z)}{z} \, dz$.\\

        % // TODO: REPASAR
        Estudiamos en primer lugar la convergencia uniforme. Dado $A\subset \Omega$, probaremos que la serie converge uniformemente en $A$ si y solo si:
        \[
            \sup\{|\Im z| : z\in A\} < \ln(3).
        \]
        \begin{description}
            \item[$\Longrightarrow)$] Demostraremos el recíproco. Como $A\subset \Omega$, no puede darse que dicho supremo sea mayor que $\ln(3)$. Supongamos por tanto que dicho supremo es $3$. Entonces, por la caracterización del supremo mediante sucesiones, existe una sucesión $\left\{\wt{z}_n\right\}$ de puntos de $A$ tal que $\left\{|\Im \wt{z}_n|\right\}\to \ln 3$. Tomando una parcial suya, existe una sucesión $\left\{z_n\right\}$ de puntos de $A$ tal que $\left|\Im z_n\right|>\ln 3 - \nicefrac{1}{n}$ para todo $n\in \bb{N}$.\\
            
            Para cada $n\in \bb{N}$, tenemos que:
            \begin{align*}
                \left|\dfrac{\cos(nz_n)}{3^n}\right| &\geq \dfrac{\left|\senh(n\Im z_n)\right|}{3^{n}}
                = \dfrac{\senh(n|\Im z_n|)}{3^{n}}
                = \dfrac{e^{n|\Im z_n|}-e^{-n|\Im z_n|}}{3^{n+1}} \geq \dfrac{e^{n|\Im z_n|}}{3^{n+1}}
                > \\
                &> \dfrac{e^{n\ln 3 -1}}{3^{n+1}}
                = \dfrac{1}{3e}
            \end{align*}

            Por tanto, como hemos encontrado una sucesión de puntos de $A$ tal que el término general no converge a $0$, entonces el término general no converge uniformemente a $0$. Por tanto, la serie no converge uniformemente en $A$.

            \item[$\Longleftarrow)$] Sea dicho supremo $r<\ln 3$. Entonces, para todo $z\in A$ y todo $n\in \bb{N}$, tenemos que:
            \begin{align*}
                \left|\dfrac{\cos(nz)}{3^n}\right| &<\left(\dfrac{e^{|\Im z|}}{3}\right)^n
                < \left(\dfrac{e^{r}}{3}\right)^n
            \end{align*}

            Como $r<\ln(3)$, tenemos que $e^{r}<3$, por lo que la base de la potencia es menor que 1. Por tanto, la serie geométrica de razón dicho cociente es convergente. Por el Test de Weierstrass, la serie de partida converge uniformemente en $A$.
        \end{description}

        En particular, dado $K\subset \Omega$ compacto, en particular es cerrado, y el supremo en cuestión se alcanza. Como es un subconjunto de $\Omega$, entonces:
        \begin{align*}
            \sup\{|\Im w| : w\in K\}=\max\{|\Im w| : w\in K\} < \ln(2)
        \end{align*}
        Por tanto, la serie converge uniformemente en compactos de $\Omega$.\\

        Para la convergencia absoluta, consideramos $z\in \Omega$ fijo. Como $\{z\}$ es compacto, por el Test de Weierstrass tenemos que converge absolutamente en $\{z\}$. Como $z$ es arbitrario, tenemos que la serie converge absolutamente en todo punto de $\Omega$.\\

        Para probar que $g$ es continua, consideramos $z\in \Omega$. Como $\Omega$ es abierto, existe $R\in \bb{R}^+$ tal que $D(z,R)\subset \Omega$. Considerando $\nicefrac{R}{2}$, tenemos que $z\in \ol{D}(z,\nicefrac{R}{2})\subset \Omega$. Por ser este compacto, la serie converge uniformemente en $\ol{D}(z,\nicefrac{R}{2})$. Como el término general de la serie es continuo para todo $n\in \bb{N}$, entonces $g$ es continua en $z$. Como $z$ es arbitrario, $g$ es continua en $\Omega$.\\

        Estudiemos ahora la convergencia uniforme del integrando en $C(0,1)$. Como $C(0,1)$ es compacto, entonces $g$ converge uniformemente y $z\mapsto \nicefrac{1}{z}$ es acotado. Por tanto, el integrando converge uniformemente en $C(0,1)$, luego podemos intercambiar la sumatoria y la integral, algo que nos será de utilidad a la hora de calcular la integral:
        \begin{align*}
            \int_{C(0,1)} \frac{g(z)}{z}\ dz &= \sum_{n=0}^{\infty} \int_{C(0,1)} \frac{\cos(nz)}{z3^n}\ dz
            = \sum_{n=0}^{\infty} \frac{1}{3^n}\int_{C(0,1)} \frac{\cos(nz)}{z}\ dz \AstIg\\&\AstIg \sum_{n=0}^{\infty} \frac{1}{3^n}\cdot 2\pi i \cdot \cos (n\cdot 0) = 2\pi i\sum_{n=0}^{\infty} \left(\dfrac{1}{3}\right)^n = 2\pi i \cdot \dfrac{1}{1-\nicefrac{1}{3}} = 3\pi i
        \end{align*}
        donde en $(\ast)$ hemos aplicado la Fórmula de Cauchy para la circunferencia, usando que el seno es una función entera.
    \end{ejercicio}

    \begin{ejercicio}[3 puntos]
        Estudiar la derivabilidad de las funciones $f, g : \mathbb{C} \to \mathbb{C}$ dadas por
        \[
            f(z) = \sen\left(\ol{z}\right)\qquad g(z) = z(z-1)f(z) \quad (z \in \mathbb{C}).
        \]

        Estudiamos en primer lugar la función $f$. Para ello, con vistas a aplicar las ecuaciones de Cauchy-Riemann, definimos $u,v:\mathbb{R}^2\to \mathbb{R}$ como:
        \begin{align*}
            u(x,y) &= \Re(f(x+iy)) = \Re(\sen(x-iy))=\sen(x)\cosh(-y)\\
            v(x,y) &= \Im(f(x+iy)) = \Im(\sen(x-iy))=\cos(x)\senh(-y).
        \end{align*}

        Calculamos ahora las derivadas parciales de $u$ y $v$:
        \begin{align*}
            \frac{\partial u}{\partial x}(x,y) &= \cos(x)\cosh(-y)\\
            \frac{\partial u}{\partial y}(x,y) &= -\sen(x)\senh(-y)\\
            \frac{\partial v}{\partial x}(x,y) &= -\sen(x)\senh(-y)\\
            \frac{\partial v}{\partial y}(x,y) &= -\cos(x)\cosh(-y)
        \end{align*}

        La primera ecuación de Cauchy-Riemann nos dice que:
        \begin{align*}
            \frac{\partial u}{\partial x}(x,y) &= \frac{\partial v}{\partial y}(x,y)\\
            \cos(x)\cosh(-y) &= -\cos(x)\cosh(-y)\\
            \cos(x)&=0\\
            x &\in  \frac{\pi}{2} + \pi\bb{Z}
        \end{align*}

        La segunda ecuación de Cauchy-Riemann nos dice que:
        \begin{align*}
            \frac{\partial u}{\partial y}(x,y) &= -\frac{\partial v}{\partial x}(x,y)\\
            -\sen(x)\senh(-y) &= \sen(x)\senh(-y)\\
            x\in \pi\bb{Z} &\lor y= 0
        \end{align*}

        Definimos por tanto el conjunto en el que se cumplen las ecuaciones de Cauchy-Riemann como:
        \begin{equation*}
            U = \left\{ z\in \bb{C} \mid \Im(z) = 0 \land \Re(z) \in \nicefrac{\pi}{2} + \pi\bb{Z} \right\}
        \end{equation*}

        Tenemos por tanto que $f\in \cc{H}(U)$, y no es derivable en ningún otro punto de $\bb{C}$.\\

        Estudiamos ahora la función $g$. Fijado $z\in \bb{C}$, veamos si es derivable en $z$.
        \begin{itemize}
            \item Si $z\in U$, entonces $g$ es derivable en $z$ por serlo $f$ y $z(z-1)$, que son funciones holomorfas en $\bb{C}$.
            \item Si $z\notin U$, distinguimos de nuevo:
            \begin{itemize}
                \item Si $z\neq 0,1$, tenemos que:
                \begin{equation*}
                    f(z) = \dfrac{g(z)}{z(z-1)}\qquad \forall z\notin U
                \end{equation*}

                Si $g$ fuese derivable en $z$, entonces $f$ también lo sería, pero sabemos que no lo es puesto que $z\notin U$. Por tanto, $g$ no es derivable en $z$.

                \item Si $z=0$, entonces:
                \begin{align*}
                    g'(0) &= \lim_{z\to 0} \frac{g(z)-g(0)}{z-0} = \lim_{z\to 0} \frac{g(z)}{z} = \lim_{z\to 0} (z-1)f(z) = -f(0) = -\sen(0) = 0
                \end{align*}

                \item Si $z=1$, entonces:
                \begin{align*}
                    g'(1) &= \lim_{z\to 1} \frac{g(z)-g(1)}{z-1} = \lim_{z\to 1} \frac{g(z)}{z-1} = \lim_{z\to 1} zf(z) = f(1) = \sen(1)
                \end{align*}
            \end{itemize}
        \end{itemize}

        Por tanto, $g$ es derivable en $U\cup \{0,1\}$ y no es derivable en ningún otro punto de $\bb{C}$.
    \end{ejercicio}

    \begin{ejercicio}[3 puntos]
        Sean $a, b \in \mathbb{C}$ con $a \neq b$ y sea $R > 0$ de modo que $R > \max\{ |a|, |b| \}$. Probar que, si $f$ es una función entera, se tiene que:
        \[
            \int_{C(0,R)} \frac{f(z)}{(z-a)(z-b)} \, dz = 2\pi i \cdot \frac{f(b) - f(a)}{b-a}.
        \]
        Deducir que toda función entera y acotada es constante (Teorema de Liouville).\\

        Descomponemos el integrando en fracciones simples:
    \begin{equation*}
        \frac{1}{(z-a)(z-b)} = \frac{A}{z-a} + \frac{B}{z-b} = \frac{A(z-b)+B(z-a)}{(z-a)(z-b)}
    \end{equation*}
    \begin{itemize}
        \item Para $z=a$: $1=A(a-b)\Longrightarrow A=\frac{1}{a-b}$.
        \item Para $z=b$: $1=B(b-a)\Longrightarrow B=\frac{1}{b-a}=-\frac{1}{a-b}$.
    \end{itemize}

    Por tanto, la integral queda:
    \begin{align*}
        \int_{C(0,R)} \frac{f(z)}{(z-a)(z-b)}dz &= \frac{1}{a-b}\left(\int_{C(0,R)} \frac{f(z)}{z-a}dz - \int_{C(0,R)} \frac{f(z)}{z-b}dz\right)\\
        &\AstIg \frac{1}{a-b}\left(2\pi i f(a) - 2\pi i f(b)\right)\\
        &= 2\pi i \cdot \frac{f(b)-f(a)}{b-a}
    \end{align*}
    donde $(\ast)$ se debe a que la función $f(z)$ es entera y que $a,b\in D(0,R)$, por lo que se puede aplicar la Fórmula de Cauchy para la circunferencia considerando como función $f(z)$.\\

    Sea ahora $f$ entera y acotada. Entonces, $\exists M\in \bb{R}^+$ tal que $|f(z)|\leq M$ para todo $z\in \bb{C}$. Por tanto, se tiene que:
    \begin{align*}
        \left|\dfrac{f(z)}{(z-a)(z-b)}\right| &\leq \frac{M}{|z-a||z-b|}
        \leq \frac{M}{\left||z| - |a|\right|\left||z| - |b|\right|}
        \leq \frac{M}{\left|R - |a|\right|\left|R - |b|\right|}
        \leq\\&\leq  \frac{M}{(R-|a|)(R-|b|)}\qquad \forall z\in C(0,R)^*
    \end{align*}
    donde hemos usado que $z\in C(0,R)^*$, por lo que $|z|=R$; y que $R>\max\{|a|,|b|\}$, por lo que $R-|a|>0$ y $R-|b|>0$. Por tanto, se tiene que:
    \begin{align*}
        \left|\int_{C(0,R)} \frac{f(z)}{(z-a)(z-b)}dz\right| &\leq \int_{C(0,R)} \left|\frac{f(z)}{(z-a)(z-b)}\right|dz\\
        &\leq 2\pi R \cdot \frac{M}{(R-|a|)(R-|b|)} = \frac{2\pi M R}{(R-|a|)(R-|b|)}
    \end{align*}

    Como la anterior expresión es válida para todo $R\in \bb{R}^+$ tal que $R>\max\{|a|,|b|\}$, podemos hacer tender $R\to \infty$. Por el Lema del Sándwich, se tiene que:
    \begin{equation*}
        \lim_{R\to \infty} \int_{C(0,R)} \frac{f(z)}{(z-a)(z-b)}dz = 0
    \end{equation*}

    Por la expresión anterior a la que habíamos llegado, se tiene que:
    \begin{equation*}
        \lim_{R\to \infty} \int_{C(0,R)} \frac{f(z)}{(z-a)(z-b)}dz = \lim_{R\to \infty} 2\pi i \cdot \frac{f(b)-f(a)}{b-a} = 2\pi i \cdot \frac{f(b)-f(a)}{b-a}
    \end{equation*}

    Por la unicidad del límite, se tiene que:
    \begin{equation*}
        f(b)=f(a) \qquad \forall a,b\in \bb{C}
    \end{equation*}

    Por tanto, $f$ es constante.
    \end{ejercicio}
\end{document}